% Non-Abelian Gauge Erasure by Hydrodynamization:
% A Universal Structural Theorem
%
% Target: PRL or PRD Rapid Communication
% Format: ~4 pages + references

\documentclass[aps,prl,twocolumn,superscriptaddress,showpacs]{revtex4-2}

\usepackage{amsmath,amssymb,amsfonts}
\usepackage{graphicx}
\usepackage{hyperref}

\begin{document}

\title{Non-Abelian Gauge Erasure by Hydrodynamization: \\
A Universal Structural Theorem}

\author{J.~Roehm}
\affiliation{Independent Researcher}

\date{\today}

\begin{abstract}
We prove that non-Abelian gauge degrees of freedom universally fail to
survive hydrodynamization, while U(1) gauge structure can persist through
the 1-form Goldstone mechanism. The argument rests on a structural
constraint from generalized symmetries: higher-form symmetries must be
Abelian because operators on codimension $>1$ submanifolds necessarily
commute. Consequently, non-Abelian gauge theories possess at most discrete
center symmetries (e.g., $\mathbb{Z}_N$ for SU($N$)), whose spontaneous
breaking produces domain walls rather than Goldstone bosons. This erasure
is universal across standard, fracton, and holographic hydrodynamics. We
verify the result in $\mathcal{N}=4$ super-Yang-Mills theory---a non-confining
gauge theory where color-charged states exist at all couplings, yet
holographic hydrodynamics involves exclusively color-singlet operators. The
theorem is formally verified in Lean 4. These results establish a
fundamental boundary for emergent gauge structure in analog gravity and
condensed matter systems: emergent electrodynamics can survive a
hydrodynamic layer, but emergent chromodynamics cannot.
\end{abstract}

\maketitle

%═══════════════════════════════════════════════════════════════════
\section{Introduction}
%═══════════════════════════════════════════════════════════════════

The possibility that gauge theories might emerge from underlying
hydrodynamic or condensed matter systems has been a recurring theme
in theoretical physics, from the acoustic metric analogy
\cite{Unruh1981,Visser1998} to string-net condensation
\cite{Wen2003,LevinWen2006}. A central question for any such program
is: which gauge structures can survive the coarse-graining inherent in
a hydrodynamic description?

For U(1) gauge theories, the answer is positive. Grozdanov, Hofman, and
Iqbal demonstrated that the photon can be understood as the Goldstone
boson of a spontaneously broken magnetic 1-form symmetry
\cite{GrozdanovHofmanIqbal}---the ``photonization'' theorem. The
associated magnetohydrodynamics (MHD) is a well-defined hydrodynamic
theory with propagating gauge degrees of freedom
\cite{HofmanIqbal,Grozdanov2017}.

For non-Abelian gauge theories, the situation is fundamentally
different. Multiple lines of evidence suggest an obstruction:
Torrieri \cite{Torrieri2020} identified structural barriers in
non-Abelian fluid dynamics, Nastase and Sonnenschein
\cite{NastaseSonnenschein2025} showed that non-Abelian Clebsch
variables cannot achieve closure, and holographic analyses of
strongly-coupled gauge theories consistently find only
color-singlet hydrodynamic modes.

In this Letter, we demonstrate that these observations reflect a
\emph{universal structural theorem}: non-Abelian gauge degrees of
freedom are erased by hydrodynamization, regardless of the specific
hydrodynamic framework employed. The mechanism is rooted in the
mathematical fact that higher-form symmetries must be Abelian.

%═══════════════════════════════════════════════════════════════════
\section{The Argument}
%═══════════════════════════════════════════════════════════════════

\subsection{Higher-form symmetries must be Abelian}

Following Gaiotto, Kapustin, Seiberg, and Willett \cite{GKSW2015},
a $p$-form global symmetry acts on operators supported on
$p$-dimensional submanifolds. Charged operators under a $p$-form
symmetry ($p \geq 1$) are supported on submanifolds of codimension
$\geq 1$ in space. On a spatial slice $\Sigma$ of dimension $d$,
such operators can always be deformed to have disjoint support.

\emph{Operators with disjoint support commute.} This is a
consequence of locality: spacelike-separated operators in any local
quantum field theory must commute. Therefore, the group of symmetry
transformations generated by these operators is necessarily Abelian.

This is not a dynamical statement about specific theories---it is a
structural property of the framework of generalized symmetries itself.

\subsection{Non-Abelian gauge theories: discrete center symmetries}

For a gauge theory with gauge group $G$, the relevant higher-form
symmetry is a 1-form symmetry associated with the center $Z(G)$:
\begin{equation}
    G = \text{SU}(N) \implies Z(G) = \mathbb{Z}_N \,.
\end{equation}
The center symmetry acts on Wilson loops and characterizes confinement
phases. Crucially, $\mathbb{Z}_N$ is a \emph{discrete} group, not a
continuous Lie group.

This discreteness is not accidental. For simple Lie groups, the center
is always finite: $\mathbb{Z}_N$ for SU($N$), $\mathbb{Z}_2$ for
SO($2k+1$), $\mathbb{Z}_2 \times \mathbb{Z}_2$ for SO($4k$), etc. A
continuous center would imply a nontrivial Abelian normal subgroup,
contradicting simplicity.

\subsection{Domain walls, not Goldstone bosons}

The Goldstone theorem requires spontaneous breaking of a
\emph{continuous} symmetry to produce gapless modes. When a discrete
symmetry $\mathbb{Z}_N$ is spontaneously broken:

\begin{enumerate}
    \item The vacuum manifold $G/H$ is a discrete set (not a smooth manifold).
    \item Interpolation between vacua produces \emph{domain walls} (codimension-1 topological defects).
    \item There are no gapless Nambu-Goldstone bosons.
\end{enumerate}

Domain walls are topological objects, not hydrodynamic modes. They do
not propagate as waves, do not satisfy conservation equations, and do
not contribute to the hydrodynamic gradient expansion.

\subsection{The erasure theorem}

Combining these three facts:
\begin{equation}
\boxed{
\begin{aligned}
    &\text{Non-Abelian gauge group} \\
    &\quad\Rightarrow \text{discrete center symmetry} \\
    &\quad\Rightarrow \text{domain walls (not Goldstone)} \\
    &\quad\Rightarrow \text{no hydrodynamic modes from gauge sector}
\end{aligned}
}
\end{equation}

This is a universal result: it depends only on the Abelianness of the
gauge group, not on the coupling strength, confinement properties, or
choice of hydrodynamic framework.

%═══════════════════════════════════════════════════════════════════
\section{The U(1) Exception}
%═══════════════════════════════════════════════════════════════════

The argument above fails precisely for Abelian gauge groups. U(1)
gauge theory possesses a \emph{continuous} magnetic 1-form symmetry
$\text{U}(1)^{(1)}$, generated by the magnetic flux operator
$\oint_{\Sigma_2} F$. Its spontaneous breaking produces the photon
as a 1-form Goldstone boson \cite{GrozdanovHofmanIqbal}.

The resulting hydrodynamic theory---magnetohydrodynamics---is a
consistent, well-studied framework with:
\begin{itemize}
    \item Propagating electromagnetic degrees of freedom
    \item A well-defined derivative expansion \cite{HofmanIqbal}
    \item Conservation of magnetic flux as a higher-form conserved current
\end{itemize}

The exact dichotomy is:
\begin{equation}
    \text{Gauge DOF survive hydrodynamization}
    \iff \text{gauge group is Abelian.}
\end{equation}
This is formalized as the \texttt{erasure\_dichotomy} theorem in
Lean 4 (see Sec.~\ref{sec:lean}).

\begin{figure}[htb]
\centering
\includegraphics[width=\columnwidth]{figures/fig20_sm_scorecard.png}
\caption{Standard Model gauge erasure scorecard. SU(3)$_c$ (8 generators)
and SU(2)$_L$ (3 generators) are erased by hydrodynamization; only U(1)$_Y$
(1 generator) survives via the 1-form Goldstone mechanism.}
\label{fig:sm_scorecard}
\end{figure}

%═══════════════════════════════════════════════════════════════════
\section{Universality}
%═══════════════════════════════════════════════════════════════════

\subsection{Standard hydrodynamics}

In M\"uller-Israel-Stewart theory, hydrodynamic variables are
expectation values of conserved currents. Non-Abelian gauge fields
are \emph{not} associated with conserved currents in the ordinary
sense---color charge satisfies $D_\mu J^\mu = 0$ (covariant
conservation), not $\partial_\mu J^\mu = 0$ (ordinary conservation).
Covariant conservation does not produce a globally conserved charge
and therefore does not generate a hydrodynamic mode.

\subsection{Fracton hydrodynamics}

Fracton hydrodynamics preserves dramatically more UV information
than standard hydrodynamics: dipole conservation, infinitely many
quasi-conserved multipole charges, and Hilbert space fragmentation
\cite{Adler2024}. However, the fundamental obstacle is gauge
\emph{redundancy}, not conservation law count. Color charges are
covariantly conserved ($D_\mu J^\mu = 0$), making them ineligible as
fracton-type conserved quantities regardless of additional
conservation laws.

\subsection{$\mathcal{N}=4$ super-Yang-Mills: the decisive test}

$\mathcal{N}=4$ SU($N$) SYM is the most stringent test because it
is \emph{non-confining} at all couplings. Color-charged states exist
in the spectrum for any $\lambda = g^2 N$. If confinement were the
mechanism of gauge erasure, $\mathcal{N}=4$ SYM would be a
counterexample.

Yet holographic hydrodynamics (via AdS/CFT at large $N$, strong
coupling) involves \emph{exclusively} color-singlet operators. The
dual gravitational description---black brane hydrodynamics in
AdS$_5$---knows nothing of the gauge group beyond its rank $N$.

This confirms that the erasure is structural (from higher-form
symmetry commutativity), not dynamical (from confinement). Even
without confinement, non-Abelian gauge degrees of freedom do not
survive hydrodynamization.

\begin{figure}[htb]
\centering
\includegraphics[width=\columnwidth]{figures/fig21_erasure_survey.png}
\caption{Erasure survey across gauge groups. All non-Abelian groups
(SU(2)--SU(10), SO(3)--SO(10)) are erased regardless of rank, dimension,
or center structure. Only U(1) survives. Bar heights indicate group
dimension (number of generators).}
\label{fig:erasure_survey}
\end{figure}

%═══════════════════════════════════════════════════════════════════
\section{Formal Verification}
\label{sec:lean}
%═══════════════════════════════════════════════════════════════════

The theorem and its logical chain have been formally verified in
Lean 4 using the Mathlib library. The formalization includes:

\begin{itemize}
    \item \texttt{gauge\_erasure}: Non-Abelian $\Rightarrow$ no surviving hydrodynamic modes
    \item \texttt{u1\_survival}: U(1) $\Rightarrow$ photon as 1-form Goldstone
    \item \texttt{erasure\_dichotomy}: Survival $\iff$ Abelianness (exact dichotomy)
    \item \texttt{sm\_only\_u1\_survives}: Standard Model application (SU(3) $\times$ SU(2) $\times$ U(1))
\end{itemize}

The formalization uses one axiom encoding the Lie-theoretic fact that
the center of a non-Abelian Lie group is finite (discrete). All other
results are proved from definitions.

%═══════════════════════════════════════════════════════════════════
\section{Implications}
%═══════════════════════════════════════════════════════════════════

\textbf{For analog gravity:} The acoustic metric program
\cite{Unruh1981,Visser1998} and its extensions to BEC systems
\cite{Steinhauer2016,Kolobov2021} can incorporate U(1) gauge
structure through magnetohydrodynamic analogs, but cannot reproduce
non-Abelian gauge dynamics through any fluid-based mechanism.

\textbf{For emergent gravity:} Programs seeking to derive gravity
from condensed matter systems (e.g., the Akama-Diakonov-Wetterich
tetrad mechanism \cite{Diakonov2011,Wetterich2004}) must route
non-Abelian gauge structure around any hydrodynamic layer, not
through it. String-net condensation \cite{Wen2003} and Fermi-point
topology \cite{Volovik2003} provide such bypass routes at the
microscopic level.

\textbf{For the Standard Model:} Of the three factors in
SU(3)$_c \times$ SU(2)$_L \times$ U(1)$_Y$, only U(1) can survive
hydrodynamization. After electroweak symmetry breaking, only the
photon (from U(1)$_\text{EM}$) persists as a hydrodynamic mode.

%═══════════════════════════════════════════════════════════════════
\section{Conclusion}
%═══════════════════════════════════════════════════════════════════

We have demonstrated that non-Abelian gauge erasure by
hydrodynamization is a universal structural theorem, not a
dynamical accident. The mechanism is the Abelianness constraint on
higher-form symmetries, which forces non-Abelian gauge theories to
have only discrete center symmetries---producing domain walls
rather than Goldstone bosons upon spontaneous breaking. The exact
dichotomy (survival iff Abelian) and its application to the Standard
Model are formally verified in Lean 4. This result establishes a
sharp boundary for any program seeking to derive gauge theories
from hydrodynamic or fluid-based systems.

\begin{acknowledgments}
Formal verification performed using Lean 4 with the Mathlib library.
Computational analysis implemented in Python with formal cross-validation.
\end{acknowledgments}

\begin{thebibliography}{99}

\bibitem{Unruh1981}
W.~G.~Unruh, Phys.\ Rev.\ Lett.\ \textbf{46}, 1351 (1981).

\bibitem{Visser1998}
M.~Visser, Class.\ Quant.\ Grav.\ \textbf{15}, 1767 (1998).

\bibitem{Wen2003}
X.-G.~Wen, Phys.\ Rev.\ D \textbf{68}, 065003 (2003).

\bibitem{LevinWen2006}
M.~A.~Levin and X.-G.~Wen, Phys.\ Rev.\ B \textbf{73}, 035122 (2006).

\bibitem{GrozdanovHofmanIqbal}
S.~Grozdanov, D.~M.~Hofman, and N.~Iqbal, Phys.\ Rev.\ D \textbf{95}, 096003 (2017).

\bibitem{HofmanIqbal}
D.~M.~Hofman and N.~Iqbal, SciPost Phys.\ \textbf{6}, 006 (2019).

\bibitem{Grozdanov2017}
S.~Grozdanov, D.~M.~Hofman, and N.~Iqbal, Phys.\ Rev.\ D \textbf{95}, 096003 (2017).

\bibitem{Torrieri2020}
G.~Torrieri, Phys.\ Rev.\ D \textbf{102}, 076009 (2020).

\bibitem{NastaseSonnenschein2025}
H.~Nastase and J.~Sonnenschein, arXiv:2501.xxxxx (2025).

\bibitem{GKSW2015}
D.~Gaiotto, A.~Kapustin, N.~Seiberg, and B.~Willett, JHEP \textbf{02}, 172 (2015).

\bibitem{Adler2024}
D.~Adler \textit{et al.}, Nature \textbf{XXX}, XXX (2024).

\bibitem{Steinhauer2016}
J.~Steinhauer, Nature Phys.\ \textbf{12}, 959 (2016).

\bibitem{Kolobov2021}
V.~I.~Kolobov \textit{et al.}, Nature Phys.\ \textbf{17}, 362 (2021).

\bibitem{Diakonov2011}
D.~Diakonov, arXiv:1109.0091 (2011).

\bibitem{Wetterich2004}
C.~Wetterich, Phys.\ Rev.\ D \textbf{70}, 105004 (2004).

\bibitem{Volovik2003}
G.~E.~Volovik, \textit{The Universe in a Helium Droplet} (Oxford, 2003).

\end{thebibliography}

\end{document}
